\chapter{Introducción}
\label{ch:introduccion}

\section{Contextualización de la Problemática}

La República Argentina representa un caso paradigmático de inestabilidad macroeconómica crónica, caracterizado por episodios recurrentes de cesación de pagos, reestructuraciones de deuda y una exposición estructural a la restricción de divisas \citep{mendoza2007, derasmo2016}. Tras la crisis de 2001 y la reestructuración de 2005, que implicó una quita sustancial de valor presente neto sobre los títulos en default, la ratio Deuda/PIB cayó de 110\% a un piso de 41,5\% hacia 2011, impulsada por el recorte nominal de capital y un ciclo de materias primas favorable. A partir de 2012 el deterioro fiscal fue continuo: el déficit primario alcanzó 3\% del PIB hacia 2016, el Estado retornó a los mercados internacionales de crédito y, tras el cierre abrupto de esos mercados en 2018 y el acuerdo con el FMI, la ratio de deuda consolidada saltó hasta 94,5\% del PIB.

La percepción de riesgo soberano, medida por el \textit{Emerging Markets Bond Index} (EMBI+), promedió 1.298 puntos básicos durante el período 2004--2025, con una desviación estándar de 1.230 puntos, lo que refleja la incompatibilidad estructural entre las condiciones de financiamiento externo y una trayectoria de deuda estable. Recién en el bienio 2024--2025 un ajuste fiscal pronunciado, que restituye el superávit primario por primera vez desde 2011, coincidió con una compresión del EMBI+ hacia 561--1.100 puntos básicos. Esta trayectoria no puede explicarse exclusivamente por decisiones de política fiscal discrecional: la hoja de balance del Banco Central de la República Argentina (BCRA), la evolución del tipo de cambio real y la restricción de divisas operaron como determinantes simultáneos, a menudo dominantes, de la ratio Deuda/PIB, y la literatura sobre \textit{fatiga fiscal} advierte que existe un límite económico e institucional para los superávits primarios que una sociedad puede sostener de manera prolongada \citep{ghosh2013}.

\section{Delimitación Científica del Objeto}

En consonancia con la ruptura epistemológica que exige toda construcción de objeto científico \citep{bachelard1938}, esta delimitación se aparta de las prenociones dicotómicas sobre la deuda pública (``buena''/``mala'', ``legítima''/``odiosa''), no operacionalizables en una matriz de datos cuantitativa. Esta tesis define la deuda pública consolidada como una restricción presupuestaria intertemporal, cuya sostenibilidad queda parametrizada por una Función de Reacción Fiscal estimable y por dinámicas estocásticas latentes en los fundamentos macroeconómicos. El riesgo soberano se modela como variable endógena que interactúa con la fatiga fiscal, y no como un factor exógeno aislado ni como resultado de la mera voluntad política del gobierno de turno.

\section{Planteamiento del Problema}

La evaluación estándar de la sostenibilidad de la deuda pública, basada en la restricción presupuestaria intertemporal de \citet{bohn1998}, sostiene que una política fiscal es sostenible si el Estado incrementa sistemáticamente su superávit primario ante el crecimiento de la ratio Deuda/PIB. Este enfoque asume que el espacio fiscal, formalizado por \citet{ostry2010} como la distancia entre el nivel de deuda vigente y el límite que la capacidad de ajuste permite sostener, es ilimitado. Dicha presunción colisiona con la realidad de las economías emergentes: modelizaciones empíricas \citep{mendoza2004, mendoza2009, ghosh2013} sugieren que la Función de Reacción Fiscal (FRF) puede exhibir no linealidades, con una reacción atenuada o inexistente bajo niveles elevados de deuda o riesgo soberano. La literatura aplicada a Argentina ha tendido además a excluir el riesgo soberano endógeno de la modelación de la reacción fiscal y a omitir el componente estocástico de las variables macroeconómicas, vacíos que esta tesis aborda directamente.

\section{Objetivo General e Hipótesis}

A partir de este planteo, la pregunta que organiza la investigación es en qué medida las trayectorias de los determinantes macroeconómicos, crecimiento del producto, tipo de cambio real, tasa de interés y resultado fiscal primario, condicionaron la solvencia intertemporal de la deuda pública consolidada argentina durante el período bajo estudio, y qué horizontes de probabilidad emergen, a partir de dicha evidencia, para el subperíodo prospectivo 2026--2035. El objetivo general de la tesis es, en consecuencia, determinar en qué medida esos determinantes macroeconómicos condicionaron la solvencia intertemporal de la deuda pública consolidada argentina, mediante la estimación de una Función de Reacción Fiscal, y proyectar su sostenibilidad hacia 2026--2035 mediante un Análisis de Sostenibilidad de la Deuda (DSA) estocástico.

La hipótesis general que vertebra la investigación sostiene que esa solvencia intertemporal estuvo condicionada por una función de reacción fiscal estructuralmente débil frente al endeudamiento, en un régimen macroeconómico caracterizado por un canal de transmisión del riesgo país (EMBI+) hacia el esfuerzo fiscal efectivo, ambos fenómenos enmarcados por la restricción externa periférica que desarrolla el Capítulo \ref{ch:marco_teorico}: no un tercer contraste independiente, sino el supuesto de fondo que encuadra la interpretación de los dos primeros.

Esa hipótesis general se desagrega en tres hipótesis derivadas, cada una vinculada a una técnica econométrica específica del protocolo del Capítulo \ref{ch:metodologia}. La primera, $H_{1a}$, sostiene que la Función de Reacción Fiscal argentina exhibe un umbral crítico de deuda o de riesgo soberano por encima del cual el coeficiente de reacción del superávit primario pierde significatividad estadística, evidenciando fatiga fiscal; se contrasta mediante el modelo de umbral de \citet{hansen1999,hansen2000} (Sección \ref{sec:hansen}). La segunda, $H_{1b}$, sostiene que el riesgo soberano no es una variable exógena respecto del esfuerzo fiscal, sino que existe una relación de doble causalidad que sesga la estimación por Mínimos Cuadrados Ordinarios simples; se contrasta mediante Variables Instrumentales en Dos Etapas (IV-2SLS). La tercera, $H_{1c}$, sostiene que la serie de la ratio Deuda/PIB exhibe al menos un quiebre estructural estadísticamente identificable en el período bajo estudio, asociado a episodios de crisis de financiamiento externo; se contrasta mediante el test de quiebre endógeno de Zivot-Andrews y el procedimiento de quiebres múltiples de Bai-Perron. Como se desarrolla en el Capítulo \ref{ch:resultados}, la evidencia recogida para estas tres hipótesis es de magnitud e intensidad estadística heterogénea, y en ningún caso alcanza, con la muestra disponible, una confirmación categórica e inequívoca.

\section{Objetivos Específicos}

Para alcanzar el objetivo general, se plantea el siguiente protocolo analítico-econométrico:

\begin{enumerate}
    \item Analizar las propiedades estocásticas de las series de tiempo macro-fiscales mediante pruebas de raíz unitaria (ADF, KPSS) y de quiebre estructural endógeno (Zivot-Andrews).
    \item Verificar la existencia de relaciones de equilibrio de largo plazo entre la ratio de endeudamiento, el esfuerzo primario, el riesgo soberano y el tipo de cambio real, mediante el contraste de Traza y Máximo Autovalor de Johansen.
    \item Estimar la Función de Reacción Fiscal histórica mediante un Modelo de Vectores con Corrección de Error (VECM) y un Vector Autorregresivo Estructural restringido (SVAR), y como ejercicio de robustez mediante Mínimos Cuadrados Ordinarios Dinámicos (DOLS).
    \item Evaluar la endogeneidad del riesgo soberano (EMBI+) respecto del esfuerzo fiscal, mediante Variables Instrumentales en Dos Etapas (IV-2SLS).
    \item Evaluar la existencia de no linealidades y regímenes de fatiga fiscal mediante un modelo de umbral univariado y su extensión multivariada, el Threshold VECM de \citet{hansen2002}.
    \item Proyectar las trayectorias de la deuda consolidada para el período 2026--2035 bajo escenarios macroeconómicos deterministas, y simular las bandas de probabilidad de un Análisis de Sostenibilidad de la Deuda estocástico mediante el método de Monte Carlo.
    \item Reconstruir y empalmar la serie trimestral de costo de financiamiento soberano a lo largo de los años de democracia ininterrumpida, identificando sus quiebres estructurales endógenos.
\end{enumerate}

\section{Justificación de la Investigación}

La justificación de esta tesis reside en el vacío de conocimiento que persiste en la literatura sobre sostenibilidad fiscal para el caso argentino, un vacío con correlato institucional directo: la Ley 27.612 de Fortalecimiento de la Sostenibilidad de la Deuda Pública \citep{ley27612} exige, desde 2021, que todo nuevo endeudamiento del Estado nacional se someta a un análisis de sostenibilidad de la deuda, sin que la literatura académica local disponga aún de un protocolo econométrico propio, replicable y estadísticamente contrastado que dé contenido empírico a esa exigencia legal. El precedente metodológico directo de ese protocolo es \citet{medeiros2012}, quien combinó por primera vez un VAR irrestricto con una Función de Reacción Fiscal de panel para simular estocásticamente la trayectoria de deuda de un conjunto de países y construir distribuciones de probabilidad en lugar de escenarios únicos, esquema luego adoptado como marco institucional por el FMI en su Marco de Riesgo Soberano y Sostenibilidad de la Deuda (SRDSF), aplicado a Argentina en 2024 (Capítulo \ref{ch:antecedentes}). Esta tesis retoma esa arquitectura general, VAR estructural más simulación estocástica, pero la especifica para el caso argentino: en la dimensión metodológica, avanza sobre las limitaciones de la ecuación de acumulación clásica en dos direcciones no exploradas conjuntamente para este período, la corrección explícita de la endogeneidad del riesgo soberano mediante instrumentación en dos etapas, y el sometimiento de la hipótesis de fatiga fiscal a un contraste de significatividad estadística formal, en lugar de asumirla por analogía con la literatura internacional. En la dimensión empírica, la consolidación de los pasivos remunerados del BCRA con la deuda del SPNF constituye una extensión que la literatura aplicada al período no ha desarrollado de manera integrada.
